% SEGMENT OLP-0496-S01
% part: intuitionistic-logic
% chapter: introduction
% section: natural-deduction

\documentclass[../../../include/open-logic-section]{subfiles}

\begin{document}

\olfileid{int}{int}{ntd}

\olsection{இயல்பான வருவித்தல்}

$\FalseCl$ விதிகள் இல்லாத இயல்பான வருவித்தல், உள்ளுணர்வுவாதத்
தருக்கத்திற்கான ஒரு நிலையான !!{derivation} அமைப்பு. விதிகளை இங்கு மீண்டும்
தந்து, BHK பொருள்கோளைப் பயன்படுத்தி அவற்றின் உந்துதலைச் சுட்டுகிறோம்.
ஒவ்வொரு நேர்விலும், முன்னுரைகளுக்கு ஆக்கங்கள் இருந்தால் முடிபுக்கும் ஓர்
ஆக்கம் உண்டு என்று முடிக்க அனுமதிக்கும் விதியாக அதைக் கருதலாம்.

இயல்பான வருவித்தல் !!{derivation}s இல் நீக்கப்படாத எடுகோள்கள் இருப்பதால்,
!!{undischarged} எடுகோள்கள்~$\Gamma$ இலிருந்து $!A$ க்கான அத்தகைய
!!a{derivation} ஐ, $!B \in \Gamma$ ஆகும் ஒவ்வொரு வாய்பாட்டின்
ஆக்கத்தையும்~$!A$ இன் ஓர்
ஆக்கமாக மாற்றும் சார்பாகக் கருத வேண்டும். !!{undischarged} எடுகோள்கள் எதுவும்
இல்லாமல்~$!A$ க்கான !!a{derivation} இருந்தால், BHK பொருள்கோளின் பொருளில்
$!A$ க்கு ஓர் ஆக்கம் உள்ளது. எனினும், விவாதத்தின் நோக்கத்திற்காகத்
தேவையில்லாதபோது~$\Gamma$ ஐக் காட்டாமல் விடுவோம்.

எடுகோள்~$!A$ தனியாகவே, !!{undischarged} எடுகோள்~$!A$ இலிருந்து~$!A$ க்கான
!!a{derivation} ஆகும். இது BHK பொருள்கோளுடன் ஒத்துப்போகிறது: ஆக்கங்கள்
மீதான முற்றொருமைச் சார்பு,~$!A$ இன் எந்த ஆக்கத்தையும்~$!A$ இன் ஓர் ஆக்கமாக
மாற்றுகிறது.

\subsection{இணைப்பு}

\begin{defish}
\AxiomC{$!A$}
\AxiomC{$!B$}
\RightLabel{\Intro{\land}}
\BinaryInfC{$!A \land !B$}
\DisplayProof
\hfill
\begin{tabular}{l}
\AxiomC{$!A \land !B$}
\RightLabel{\Elim{\land}}
\UnaryInfC{$!A$}
\DisplayProof
\\[3ex]
\AxiomC{$!A \land !B$}
\RightLabel{\Elim{\land}}
\UnaryInfC{$!B$}
\DisplayProof
\end{tabular}
\end{defish}

முறையே $!A_1$, $!A_2$ ஆகியவற்றின் ஆக்கங்கள் $N_1$, $N_2$ நம்மிடம்
இருப்பதாகக் கொள்க. அப்போது $!A_1 \land !A_2$ இன் ஆக்கம், அதாவது
$\tuple{N_1, N_2}$ என்ற சோடியும் நம்மிடம் உள்ளது.

BHK பொருள்கோளின்படி $!A_1 \land !A_2$ இன் ஆக்கம் என்பது $\tuple{N_1,
N_2}$ என்ற சோடி. அத்தகைய சோடி நம்மிடம் இருப்பதாகக் கொள்க. அப்போது ஒவ்வோர்
இணைப்புறுப்பின் ஆக்கமும் நம்மிடம் உள்ளது: $N_1$ ஆனது~$!A_1$ இன் ஆக்கம்;
$N_2$ ஆனது~$!A_2$ இன் ஆக்கம்.
% SOURCE CORRECTION TA-IL-003: the conjunction has distinct A_1 and A_2 conjuncts.

\subsection{நிபந்தனைவாக்கம்}

\begin{defish}
  \AxiomC{$\Discharge{!A}{u}$}
  \DeduceC{$!B$}
  \DischargeRule{$\Intro{\lif}$}{u}
  \UnaryInfC{$!A \lif !B$}
  \DisplayProof
\hfill
  \AxiomC{$!A \lif !B$}
  \AxiomC{$!A$}
  \RightLabel{$\Elim{\lif}$}
  \BinaryInfC{$!B$}
  \DisplayProof
\end{defish}

!!{undischarged} எடுகோள்~$!A$ இலிருந்து~$!B$ க்கான !!a{derivation} நம்மிடம்
இருந்தால்,~$!A$ இன் ஆக்கங்களை~$!B$ இன் ஆக்கங்களாக மாற்றும் சார்பு~$f$
ஒன்று உள்ளது. அதே சார்பு $!A \lif !B$ இன் ஓர் ஆக்கம். எனவே $\Intro\lif$
இன் முன்னுரைக்கு,~$!A$ இன் ஓர் ஆக்கத்தைச் சார்ந்த ஓர் ஆக்கம் இருந்தால்,
முடிவு $!A \lif !B$ க்கு ஓர் ஆக்கம் உள்ளது.

மறுபுறம், $!A$ இன் ஆக்கம்~$N$ உம் $!A \lif !B$ இன் ஆக்கம்~$f$ உம்
இருப்பதாகக் கொள்க. $!A \lif !B$ இன் ஆக்கம் என்பது~$!A$ இன் ஆக்கங்களை~$!B$
இன் ஆக்கங்களாக மாற்றும் சார்பு. எனவே $f(N)$ ஆனது~$!B$ இன் ஓர் ஆக்கம்;
அதாவது $\Elim\lif$ இன் முடிபுக்கு ஓர் ஆக்கம் உள்ளது.

\subsection{பிரிப்பு}

\begin{defish}
\begin{tabular}{l}
\AxiomC{$!A$}
\RightLabel{\Intro{\lor}}
\UnaryInfC{$!A \lor !B$}
\DisplayProof
\\[3ex]
\AxiomC{$!B$}
\RightLabel{\Intro{\lor}}
\UnaryInfC{$!A \lor !B$}
\DisplayProof
\end{tabular}
\hfill
\AxiomC{$!A \lor !B$}
\AxiomC{$\Discharge{!A}{n}$}
\DeduceC{$!C$}
\AxiomC{$\Discharge{!B}{n}$}
\DeduceC{$!C$}
\DischargeRule{\Elim{\lor}}{n}
\TrinaryInfC{$!C$}
\DisplayProof
\end{defish}

$!A_i$ இன் ஆக்கம்~$N_i$ நம்மிடம் இருந்தால், அதை~$!A_1 \lor !A_2$ இன்
ஆக்கம்~$\tuple{i, N_i}$ ஆக மாற்றலாம். மறுபுறம், $!A_1 \lor !A_2$ இன் ஓர்
ஆக்கம், அதாவது $N_i$ ஆனது~$!A_i$ இன் ஆக்கமாக இருக்கும் $\tuple{i, N_i}$
என்ற சோடி நம்மிடம் இருப்பதாகவும், முறையே $!A_1$, $!A_2$ ஆகியவற்றின்
ஆக்கங்களை~$!C$ இன் ஆக்கங்களாக மாற்றும் சார்புகள் $f_1$, $f_2$ நம்மிடம்
இருப்பதாகவும் கொள்க. அப்போது $f_i(N_i)$ ஆனது~$!C$ இன் ஓர் ஆக்கம்; இதுவே
$\Elim\lor$ இன் முடிபாகும்.

\subsection{பொருத்தமின்மை}

\begin{defish}
  \AxiomC{$\lfalse$}
  \RightLabel{\FalseInt}
  \UnaryInfC{$!A$}
  \DisplayProof
\end{defish}

!!{undischarged} எடுகோள்கள்~$!B_1$, \dots, $!B_n$ இலிருந்து~$\lfalse$
க்கான !!a{derivation} நம்மிடம் இருந்தால், $!B_1$, \dots, $!B_n$
ஆகியவற்றின் ஆக்கங்களை~$\lfalse$ இன் ஓர் ஆக்கமாக மாற்றும் சார்பு
$f(M_1, \dots, M_n)$ ஒன்று உள்ளது. $\lfalse$ க்கு ஆக்கம் இல்லாததால்,
$!B_1$, \dots, $!B_n$ ஆகிய அனைத்துக்கும் ஆக்கங்கள் இருக்க முடியாது.
எனவே $f$ க்குப் பின்வரும் பண்பும் உண்டு: $M_1$, \dots, $M_n$ ஆகியவை முறையே
$!B_1$, \dots, $!B_n$ ஆகியவற்றின் ஆக்கங்களாக \emph{இருந்தால்}, $f(M_1,
\dots, M_n)$ ஆனது~$!A$ இன் ஓர் ஆக்கமாக \emph{இருக்கும்}.

\subsection{$\lnot$ க்கான விதிகள்}

$\lnot !A$ ஆனது $!A \lif \lfalse$ என வரையறுக்கப்படுவதால், கண்டிப்பாகச்
சொன்னால்~$\lnot$ க்கான விதிகள் நமக்குத் தேவையில்லை. ஆனால் அவற்றைத்
தந்தால், அவை பின்வருமாறு இருக்கும்:

\begin{defish}
\AxiomC{$\Discharge{!A}{n}$}
\noLine
\DeduceC{$\lfalse$}
\DischargeRule{\Intro{\lnot}}{n}
\UnaryInfC{$\lnot !A$}
\DisplayProof
\hfill
\AxiomC{$\lnot !A$}
\AxiomC{$!A$}
\RightLabel{\Elim{\lnot}}
\BinaryInfC{$\lfalse$}
\DisplayProof
\end{defish}


\subsection{\usetoken{P}{derivation} எடுத்துக்காட்டுகள்}

\begin{enumerate}
\item $\Proves !A \lif (\lnot !A \lif \lfalse)$, அதாவது $\Proves !A
  \lif ((!A \lif \lfalse) \lif \lfalse)$
  \begin{prooftree}
    \AxiomC{$\Discharge{!A}{2}$}
    \AxiomC{$\Discharge{!A \lif \lfalse}{1}$}
    \RightLabel{\Elim\lif}
    \BinaryInfC{$\lfalse$}
    \DischargeRule{\Intro\lif}{1}
    \UnaryInfC{$(!A \lif \lfalse)\lif \lfalse$}
    \DischargeRule{\Intro\lif}{2}
    \UnaryInfC{$!A \lif (!A \lif \lfalse) \lif \lfalse$}
  \end{prooftree}

\item $\Proves ((!A \land !B) \lif !C) \lif (!A \lif (!B \lif !C))$
  \begin{prooftree}
    \AxiomC{$\Discharge{(!A \land !B) \lif !C}{3}$}
    \AxiomC{$\Discharge{!A}{2}$}
    \AxiomC{$\Discharge{!B}{1}$}
    \RightLabel{\Intro\land}
    \BinaryInfC{$!A \land !B$}
    \RightLabel{\Elim\lif}
    \BinaryInfC{$!C$}
    \DischargeRule{\Intro\lif}{1}
    \UnaryInfC{$!B \lif !C$}
    \DischargeRule{\Intro\lif}{2}
    \UnaryInfC{$!A \lif (!B \lif !C)$}
    \DischargeRule{\Intro\lif}{3}
    \UnaryInfC{$((!A \land !B) \lif !C) \lif (!A \lif (!B \lif !C))$}
  \end{prooftree}

\item $\Proves \lnot(!A \land \lnot !A)$, அதாவது $\Proves (!A \land (!A
  \lif \lfalse))\lif \lfalse$
  \begin{prooftree}
    \AxiomC{$\Discharge{!A \land (!A \lif \lfalse)}{1}$}
    \RightLabel{\Elim\land}
    \UnaryInfC{$!A \lif \lfalse$}
    \AxiomC{$\Discharge{!A \land (!A \lif \lfalse)}{1}$}
    \RightLabel{\Elim\land}
    \UnaryInfC{$!A$}
    \RightLabel{\Elim\lif}
    \BinaryInfC{$\lfalse$}
    \DischargeRule{\Intro\lif}{1}
    \UnaryInfC{$(!A \land (!A \lif \lfalse))\lif \lfalse$}
  \end{prooftree}

\item $\Proves \lnot\lnot(!A \lor \lnot !A)$, அதாவது $\Proves ((!A \lor
  (!A \lif \lfalse))\lif \lfalse)\lif \lfalse$
  \begin{prooftree}\hskip-3em
    \AxiomC{$\Discharge{(!A \lor (!A \lif \lfalse))\lif \lfalse}{2}$}
    \AxiomC{$\Discharge{(!A \lor (!A \lif \lfalse))\lif \lfalse}{2}$}
    \AxiomC{$\Discharge{!A}{1}$}
    \RightLabel{\Intro\lor}
    \UnaryInfC{$!A \lor (!A \lif \lfalse)$}
    \RightLabel{\Elim\lif}
    \BinaryInfC{$\lfalse$}
    \DischargeRule{\Intro\lif}{1}
    \UnaryInfC{$!A \lif \lfalse$}
    \RightLabel{\Intro\lor}
    \UnaryInfC{$!A \lor (!A \lif \lfalse)$}
    \RightLabel{\Elim\lif}
    \insertBetweenHyps{\hskip -4em}
    \BinaryInfC{$\lfalse$}
    \DischargeRule{\Intro\lif}{2}
    \UnaryInfC{$((!A \lor (!A \lif \lfalse))\lif \lfalse)\lif \lfalse$}
  \end{prooftree}
\end{enumerate}

\begin{prop}
  உள்ளுணர்வுவாதத் தருக்கத்திற்கான இயல்பான வருவித்தலில் $\Gamma \Proves !A$
  எனில், செவ்வியல் தருக்கத்திலும் $\Gamma \Proves !A$. குறிப்பாக, $!A$ ஓர்
  உள்ளுணர்வுவாதத் தேற்றம் எனில், அது ஒரு செவ்வியல் தேற்றமும் ஆகும்.
\end{prop}

\begin{proof}
  ஒவ்வோர் இயல்பான வருவித்தல் விதியும் செவ்வியல் இயல்பான வருவித்தலிலும் ஒரு
  விதியாகும்; எனவே உள்ளுணர்வுவாதத் தருக்கத்தின் ஒவ்வோர் !!{derivation} உம்
  செவ்வியல் தருக்கத்தின் ஓர் !!a{derivation} ஆகும்.
\end{proof}

\begin{prob}
  பின்வரும் !!{formula}s க்கு உள்ளுணர்வுவாதத் தருக்கத்தில்
  !!{derivation}s தருக:
  \begin{enumerate}
    \item $(\lnot !A \lor !B) \lif (!A \lif !B)$
    \item $\lnot\lnot\lnot !A \lif \lnot !A$
    \item $\lnot\lnot (!A \land !B) \liff (\lnot\lnot !A\land \lnot \lnot !B)$
    \item $\lnot (!A \lor !B) \liff (\lnot !A \land \lnot !B)$
    \item $(\lnot !A \lor \lnot !B) \lif \lnot (!A \land !B)$
    \item $\lnot\lnot ( !A \land !B) \lif  (\lnot\lnot !A \lor
    \lnot\lnot !B)$
  \end{enumerate}
\end{prob}

\end{document}
